\chap{评价与资源分配：怎样判断系统放大了能力}
\label{sec:evaluation}
\subsection{从产出计数到工程结果}
评价智能开发系统时，很容易获取代码行数、消息数、工具调用数和 token 数。这些量描述活动，但不能独自说明工程价值。更多代码可能增加维护负担，更多报告可能消耗人的注意力，更多测试也可能重复检查已经清楚的事实。评价应首先联系任务结果、使用质量和后续变化能力。

开发者生产率研究中的 SPACE 框架强调多维度观察，避免依赖单一活动指标\cite{space}。在 Agent 开发场景中，同样需要同时观察产物、过程、协作和人的体验。本文建议保留多个维度的结果，在明确质量要求和资源条件之后比较，而不是急于合成一个普适分数。

能力放大可以在受控条件下定义为相对于基线的增益。基线是直接把相同目标交给一个具备相同工具和资源条件的强 Agent。组织方式增加之后，若交付质量提高、人类时间下降或可承担复杂度扩大，则可以进一步分析收益来源。基线也应经过合理配置，避免用弱基线制造优势。

\subsection{核心评价维度与记录口径}
\begin{table}[htbp]\centering\small
\caption{建议的评价维度与观察对象}\label{tab:metrics}
\begin{tabularx}{\linewidth}{p{2.6cm}YY}\toprule
维度 & 主要观察对象 & 解释时需要注意\\\midrule
任务结果 & 接受条件、实际行为、未完成项 & 保持独立检查，防止成功定义漂移\\
使用质量 & 完成任务、错误理解、绕路和等待 & Agent 成功不等于普通用户顺畅\\
演化能力 & 后续修改成本、影响范围、返工 & 需要多轮任务与维护观察\\
交付效率 & 需求到可用、故障到有效修复的时间 & 区分执行、等待与恢复时间\\
人类负担 & 理解、比较、表达、纠偏、恢复时间 & 同时观察理解正确性与决策质量\\
系统投入 & token、费用、工具时间和协调开销 & 记录模型、缓存与计费口径\\\bottomrule
\end{tabularx}\end{table}

对人类负担的观察可以使用任务计时、操作记录与简短理解题。比如，阅读汇报后请参与者判断取消是否已经完成、指出某项协议的责任方，或选择符合给定目标的方案。若阅读更快但理解错误更多，不能视为改进。主观满意度有价值，但需要与实际任务表现一起解释。

架构质量可以通过后续变化体现。给系统加入一项相关但未提前公开的需求，观察需要理解和修改哪些模块、是否破坏既有契约，以及是否增加特殊分支。代码结构指标可以辅助描述，真正重要的是局部变化是否迫使无关范围同步调整。

\subsection{多轮任务与比较条件}
一次功能实现不足以评价长期组织。可以设计连续任务序列：新增能力、改变约束、处理跨模块故障、完成交接，再实现另一项依赖此前经验的需求。各条件从相同初始状态开始，使用相同任务和外部评价。后续任务需要避免被某条件提前获得，以便观察真实适应能力。

比较条件可以逐步增加机制：单 Agent 基线；沿模块组织的团队；团队加持续经验；团队加智能交互；完整框架。消融实验可以去除某个机制，检查观察到的收益是否仍存在。不同条件应记录实际资源，并同时报告预算固定和质量目标固定的比较方式。

\fig{15-evaluation}{多轮演化评价设计。不同组织条件面对相同任务序列，以独立检查和人类任务共同评价结果，并记录资源与失败。}{evaluation}

人类参与实验需要考虑学习效应。参与者看过一个界面后，会对同一问题更熟悉，因此不同表达方式的顺序应平衡，或使用难度相近的不同任务。开发系统也可能因随机生成和工具波动产生差异，需要重复运行并报告分布，而不只挑选最佳案例。

模型升级后的评价应同时保留可比任务和增加更复杂任务。前者帮助了解同一工作是否变得更高效，后者帮助观察方法是否扩展能力边界。若只保留已经变得容易的任务，可能低估组织方法的潜在作用；若只增加难度，又难以区分模型与任务变化带来的影响。

\subsection{展示与交互可能成为主要投入}
本文提出一个值得观察的成本结构假设：面向人类理解和控制的 token 消耗，可能与代码生成处于相当规模。这类工作包含选择重要信息、解释背景、生成交互演示、核对状态和根据反馈更新表达。它们可能贯穿整个维护周期，尤其在软件复杂度扩大之后持续发生。

可以把模型投入按活动记为
\begin{equation}
 T=T_{\text{实现}}+T_{\text{验证}}+T_{\text{协调}}+T_{\text{表达}}+T_{\text{其他}},
\end{equation}
并观察 $T_{\text{表达}}/T_{\text{实现}}$ 随任务、用户熟悉程度和复用策略怎样变化。这个比例是描述量，不是预算目标。同一次调用若同时服务多个活动，应采用一致归类规则，或明确记录共享部分，避免重复计数。

表达投入的收益可能通过人的理解时间下降、反馈更准确和后续返工减少体现。一个生成演示需要更多 token，却帮助人提前发现关键误解，可能降低总成本。反过来，频繁生成无人查看的复杂报告，即使内容正确也可能没有足够收益。评价需要把使用和后续结果联系起来。

\fig{16-cost}{表达投入的价值假设。增加生成资源可能改善理解与反馈，但是否值得需要通过人的任务表现和后续工程结果判断；图中没有预设最佳比例。}{cost}

成本还受复用方式影响。稳定模板、原生组件和已验证示例可以降低后续生成成本；局部更新可以避免重复制作未变化内容；按需生成适合少用的深入视图。不同策略可能同时改变等待时间与用户体验，因此应把延迟和资源一起报告。

\subsection{可检验的研究假设}
第一项假设是，沿稳定职责组织持续成员能够减少后续任务的背景重建与跨阶段交接。可以观察调查时间、重复工具读取、契约冲突和维护质量。若模块边界不清，或者任务变化完全跨越既有职责，这种收益可能减弱。

第二项假设是，基于变化与证据选择验证范围，能够在保持重要缺陷发现能力的同时减少无效检查。需要同时记录漏检、误报、检查费用和修复延迟，避免只比较运行次数。对于全量测试极便宜的项目，自适应选择的管理成本可能高于收益。

第三项假设是，把理解与控制放在同一界面，能够提高反馈的准确性并降低上下文重述成本。可以比较纯文字汇报、固定仪表盘和按任务生成的交互视图，让参与者完成相同理解与决策任务。若生成视图经常改变结构，学习成本可能抵消优势。

第四项假设是，持续经验与智能交互的收益会在多轮维护中累积。验证它需要较长观察，并控制模型、工具和任务顺序的变化。一次成功展示只能说明该情境可行，不能证明长期收益已经成立。

\subsection{一个可实施的首轮研究方案}
首轮研究可以聚焦交互投影与多轮维护的组合。选择具有清楚接受条件的任务平台，准备若干难度相近的变化序列，每个序列包含实现、需求调整和故障修复。先在固定模型与工具条件下运行单 Agent 和模块团队，再在相同成果上比较文字汇报与交互投影。这样可以分别观察组织和表达的作用。

参与者接收明确目标，例如确认某项行为能否使用、找出一个故意保留的解释缺口，以及提出符合目标的下一步改进。记录完成时间、判断正确性、反馈能否被执行者准确定位，以及执行者据此修改所需的额外澄清。界面的生成时间和费用也纳入记录，才能计算整个协作过程的成本。

研究应在开始前明确主要结果和停止条件，并报告全部完成与未完成运行。样本量可依据预实验中的波动和希望识别的效应确定，本文不预设任意人数。分析时同时展示差异分布和典型失败过程，使读者能够判断收益来自哪些条件，而不只看到一个平均分数。

\subsection{哪些观察会要求修改框架}
一个可检验框架需要说明什么结果会削弱其主张。若模块团队在多轮任务中持续花费更多协调时间，且没有改善结果或维护质量，就应重新判断职责划分是否有价值。若长期记忆增加了错误假设的复用，则需要调整经验维护，而不能仅以“积累不足”解释所有负面结果。

\begin{table}[htbp]\centering\small
\caption{负面结果及其可能要求的设计调整}\label{tab:falsifiers}
\begin{tabularx}{\linewidth}{p{3.6cm}YY}\toprule
观察 & 需要进一步区分的原因 & 可能的调整\\\midrule
团队协调超过并行收益 & 边界不清或任务过小 & 合并职责、减少层级、改变委派粒度\\
生成界面阅读更快但判断更差 & 信息缺失或视觉暗示误导 & 补充关键条件，重新组织状态解释\\
自适应选测仍遗漏重要缺陷 & 依赖信息不完整或选择偏差 & 扩大检查，修订证据关联\\
经验复用导致重复错误 & 条件过时或缺少反例 & 缩小经验范围，增加核对入口\\
展示投入增长却很少被使用 & 关注点不符或生成过度 & 按需生成，减少重复更新\\\bottomrule
\end{tabularx}\end{table}

同样，若一个更强的单 Agent 在相同条件下表现更好，框架可以采用更简化的组织，同时继续使用目标、证据和交互机制。能力放大的目标要求机制接受这种调整。研究应报告失败和适用边界，使下一次设计能够选择更合适的组合。

\subsection{效度威胁与结果解释}
任务选择、评价泄漏、模型随机性和参与者经验都可能影响结论。开发者提供的示例也可能更适合某种组织。评价应公开任务条件与失败情况，保留独立接受条件，并记录因环境问题未完成的运行。环境失败与产品失败应分别描述，但不能简单删除所有不利样本。

对结果的解释需要关注适用范围。一个小型项目的组织开销可能很高，一个长期多模块系统则可能受益明显。一个熟悉系统的用户可能偏好紧凑图示，首次接触者需要更多实例。报告这些条件，比宣称某种流程普遍优越更有帮助。

本文提出的评价方案尚未在本稿中执行。后续研究可以先用少量完整任务验证记录与测量是否可行，再扩大项目和参与者范围。评价过程本身应避免成为新的高成本仪式，优先选择能区分主要假设的观察。
